\scanpage{58}

Let us also put
\[
(90)\qquad
\left\{\begin{aligned}
\nu'&=\nu_3-\nu_2-3\nu_2\nu_3+3\nu_0\\
\nu''&=\nu_3-\nu_2\unsure{\deleted{+}-|\deleted{-}+}2\nu_2\cdot\nu_3+2\nu_0.
\end{aligned}\right.
\]

We then have the relations
\[
(91)\qquad
\left\{\begin{aligned}
\mu&=1-2\nu_2+4\nu'&&({}=\varepsilon_2+4\nu')\\
\mu&=1-2\nu_3+6\nu''&&({}=\varepsilon_3+6\nu'')\\[6pt]
&\cdot\\[-3pt]
\underbrace{-\nu_2+2\nu'}_{\nu}&=
\underbrace{-\nu_3+3\nu''}_{\nu}
&&\text{i.e. }2\nu'-\nu_2+\nu_3=3\nu''.
\end{aligned}\right.
\]
\margin{where $\varepsilon_p=(-1)^{\nu_p}$
\par whence \deleted{$\unsure{\varepsilon_p=1-2\nu_p}$}
\par \deleted{whence} $\varepsilon_p-1=-2\nu_p$
\par $(p\in\{2,3\})$}

\uline{N.B.} The condition $\underline\mu=1$ is equivalent to the conditions
\[
(92)\qquad \nu_2=\nu_3=\nu_0=0
\]
and implies
\[
(93)\qquad \nu=\nu'=\nu''=0
\]

\scanpage{59}
\mspage{677}

4. \unsure{The groups|The group}
\deleted{$\widetilde Z_{\rho,\sigma}$ \unsure{and the sequence} \illegible{2}}
$N_{\rho,\sigma}$, \deleted{\unsure{exact}}\inserted{and the}
exact sequence
\diagram{\begin{tikzcd}[column sep=small]
1 \arrow[r] & N_{\rho,\sigma}^{0+} \arrow[r] & N_{\rho,\sigma} \arrow[r] & \mathbb O_{12}^{*} \arrow[r,dash] & 1.
\end{tikzcd}}
Relations with $\mathrm{Gl}(2,\mathbb Z/12\mathbb Z)$.

We define it as embedded in
$\mathbb O_{12}^{*}\times\underline M\times\underline M$:
\[
\begin{aligned}
(94)\quad&\widetilde Z_{\rho,\sigma}
\subset \mathbb O_{12}^{*}\times\underline M\times\underline M\\[-2pt]
&\prime\prime\\[-2pt]
&\bigl\{(\underline\mu,\alpha_0,\beta_0)\mid
(1-\rho_{\underline M})(\alpha_0)-(1-\sigma_{\underline M})(\beta_0)
=3\nu''(\underline\mu)\lambda_1\bigr\}
\end{aligned}
\]

Note that we have
\[
(95)\qquad (1-\rho_{\underline M})
(\underbrace{\lambda_1-\lambda_0})=3\lambda_1
\]
so the cocycle relation in (\unsure{72|77})
thus takes the form
\[
(96)\qquad
(1-\rho_{\underline M})(\alpha_0-\nu''(\lambda_1-\lambda_0))
-(1-\sigma_{\underline M})(\beta_0)=0
\]
I shall assume \unsure{again} that we have a natural group structure on
$\widetilde Z_{\rho,\sigma}$, for which the projection
$\widetilde Z_{\rho,\sigma}\to\mathbb O_{12}^{*}$
is a group homomorphism. We now find an exact sequence

(97)
\diagram{\begin{tikzcd}[column sep=large]
1 \arrow[r,dash] & \widetilde Z_{\rho,\sigma}^{+0\deleted{\unsure{*}}}
\arrow[r] & \widetilde Z_{\rho,\sigma} \arrow[r] & \mathbb O_{12}^{*} \arrow[r,dash] & 1
\end{tikzcd}}
where
\[
\begin{aligned}
(98)\quad&\widetilde Z_{\rho,\sigma}^{+0}
\overset{\mathrm{dfn}}=\operatorname{Ker}
(\widetilde Z_{\rho,\sigma}\mathbin{-}\mathbb O_{12}^{*})\\[-2pt]
&\rotatebox{90}{$\simeq$}\\[-2pt]
&\bigl\{(\alpha_0,\beta_0)\mid
(1-\rho_{\underline M})(\alpha_0)-(1-\sigma_{\underline M})(\beta_0)=0\bigr\}
\end{aligned}
\]
This group scheme is smooth in characteristic $\ne3$, but not in characteristic $3$, where its dimension goes

\scanpage{60}

from relative dimension $2$ to relative dimension $3$.
This fact of an exceptional fiber of larger dimension cannot, it seems to me, be
eliminated by a dodge, \deleted{\unsure{of}} by introducing
\deleted{some} some ad-hoc definitions for a ``better'' group
$\widetilde Z_{\rho,\sigma}$: there are ``more'' automorphisms of
$\underline{\unsure{\mathcal G}}_k$ (for $k$ an
algebraically closed field), having the properties envisaged, if $k$ is of characteristic $3$, than in the other cases of characteristic.

However, it seems to me that we can construct a subgroup scheme
\deleted{\unsure{of}} of $\widetilde Z_{\rho,\sigma}$, which I feel like denoting
$N_{\rho,\sigma}$, which is smooth over $S_0=\operatorname{Spec}(\mathbb Z)$ and
coincides with $\widetilde Z_{\rho,\sigma}$ over
\deleted{$\unsure{S_{0,3}}$}
$S_{0,\unsure{3|2}}\unsure{{}=}\operatorname{Spec}\mathbb Z[1/3]=S_0\setminus\{3\}$.
To describe it, I note that in the loop equation\inserted{ (96)},
the \deleted{right-hand}\inserted{left-hand} side is
\unsure{automatically|automatic} in
\deleted{$\underline M$, \unsure{which}}\inserted{the submodule of $\underline M$ corresponding}
to the congruence
\[
c_0+c_1\equiv0\ (3)\qquad(\mathrm{cf}\ 62)
\]
and if we introduce the corresponding sub-\deleted{group}\inserted{$\mathbb Z$-module} of index

\scanpage{61}
\mspage{678}

$3$, $M_0$ of $M$
\[
\begin{aligned}
(99)\qquad&M_0\subset M\\[-2pt]
&\prime\prime\ \mathrm{dfn}\\[-2pt]
&\{c_0\lambda_0+c_1\lambda_1\mid c_0+c_1\equiv0\ (3)\},
\end{aligned}
\]
whence a vector bundle $\underline M_0$ over $S_0$ ---
we in fact have a morphism of schemes

(100)
\diagram{\begin{tikzcd}[column sep=large]
\widetilde Z_{\rho,\sigma} \arrow[r] & \underline M_0\\
(\underline\mu,\alpha_0,\beta_0) \arrow[r,mapsto] &
(1-\rho_{\underline M})
(\underbrace{\alpha_0+\nu''(\lambda_1-\lambda_0)}_{\widetilde\alpha_0})
+(1-\sigma_{\underline M})(\beta_0)
\end{tikzcd}}
which we can make explicit as follows: we have a homomorphism of \uline{$\mathbb Z$-modules}

(101)
\diagram{\begin{tikzcd}[column sep=large]
M\times M\times\mathbb Z \arrow[r] & M_0\\
(\alpha_0,\beta_0,\nu'') \arrow[r,mapsto] &
(1-\rho_M)(\alpha_0+\nu''(\lambda_1-\lambda_0))
+(1-\sigma_M)(\unsure{\beta_0})
\end{tikzcd}}
whence a homomorphism \unsure{for} the corresponding module schemes

(102)
\diagram{\begin{tikzcd}[column sep=large]
\underline M\times\underline M\times\underbrace{\underline{\mathbb Z}}_{\mathbb G_a}
\arrow[r] & \underline{M_0}
\end{tikzcd}}
which we compose with

(103)
\diagram{\begin{tikzcd}[column sep=large]
\widetilde Z_{\rho,\sigma} \arrow[r] & \underline M\times\underline M\times\underline{\mathbb Z}\\
(\underline\mu,\alpha_0,\beta_0) \arrow[r,mapsto] & (\alpha_0,\beta_0,\nu''(\mu))
\end{tikzcd}}
We can then consider the cocycle equation (96) as an equation \uline{in}
$\underline M_0$ --- it thus becomes stricter,

\scanpage{62}

whence a closed subscheme of $\widetilde Z_{\rho,\sigma}$,
\[
\begin{aligned}
(104)\quad&N_{\rho,\sigma}\subset\widetilde Z_{\rho,\sigma}\\[-2pt]
&\prime\prime\\[-2pt]
&\{(\underline\mu,\alpha_0,\beta_0)\in\widetilde Z_{\rho,\sigma}\mid
\text{the \deleted{cocycle}\inserted{loop} equation}\\
&\hspace{5cm}\text{(96) valid in }\underline M_0\}.
\end{aligned}
\]
Let us then consider the morphism
\diagram{\begin{tikzcd}
N_{\rho,\sigma} \arrow[r] & \mathbb O_{12}^{*}
\end{tikzcd}}
\margin{N.B. The group $N_{\rho,\sigma}$ should also be denotable by
$\unsure{M|\widetilde N}_{\rho,\sigma}[0]$, unless I am mistaken it
\par \unsure{Or|In} rather it would be
$(\unsure{M}_{\rho,\sigma}^{\unsure{\sharp|b}}/\unsure{M}_{\rho})[0]$
\deleted{\unsure{according to}} \unsure{with} the notations of
$\S\,49$, \deleted{\unsure{2}}\inserted{\uline{$\mathrm{VI}$}}}
It is clear that this morphism is still an epimorphism, the fiber over a
$\underline\mu=(\nu_2,\nu_3,\nu_0)$ being made up of the
$(\alpha_0,\beta_0)\in\Gamma\,\underline M\times\underline M$
satisfying (96) in $M_0$,
\deleted{We \unsure{take the composite homomorphism}}
\inserted{\unsure{is identified|being identified} with the \unsure{scheme} of}
$(\widetilde\alpha_0,\beta_0)$ satisfying the linear equation
\diagram{\begin{tikzcd}
\deleted{M\times M\times\unsure{\mathbb Z}}
\arrow[r,phantom,"{\deleted{\longrightarrow}}" description]
& \deleted{M_0}
\end{tikzcd}}
in $\widetilde\alpha_0,\beta_0$
\[
(105)\qquad
(1-\rho_{\underline M})(\widetilde\alpha_0)
-(1-\sigma_{\underline M})(\beta_0)=0
\qquad\inserted{\text{in }\underline M_0}.
\]
Now the homomorphism of Modules

(106)
\diagram{\begin{tikzcd}[column sep=large]
{} \arrow[d,phantom,"{\Bigg\{}" description] & \underline M\times\underline M \arrow[r] & \underline M_0\\
{} & (\widetilde\alpha_0,\beta_0) \arrow[r,mapsto] &
(1-\rho_{\underline M})(\widetilde\alpha_0)-(1-\sigma_{\underline M})(\beta_0)
\end{tikzcd}}
is a smooth epimorphism --- being associated to a corresponding homomorphism
\diagram{\begin{tikzcd}
M\times M \arrow[r] & M_0
\end{tikzcd}}
which is \unsure{now|manifestly} an epimorphism.
Thus \deleted{$\unsure{\widetilde Z}$} $N_{\rho,\sigma}$ \uline{is smooth} over $S_0$.
\deleted{\unsure{This}} Being a closed subscheme of $\widetilde Z_{\rho,\sigma}$

\scanpage{63}
\mspage{679}

which coincides with it over $S_0\setminus\{3\}$, it
follows that it is the \uline{schematic closure}
\unsure{that|by}
\[
\begin{gathered}
(107)\qquad
N_{\rho,\sigma}=\widetilde Z_{\rho,\sigma}|(\mathrm{Spec}(\mathbb Z)\setminus\{3\}),\\
\text{(schematic closure)}
\end{gathered}
\]
which implies that it is also a subgroup scheme.

Perhaps this is the moment to look a little more closely at the
sub-$\mathbb Z$-module $M_0$ of $M$,
\[
\begin{aligned}
(108)\qquad M_0
&=\mathbb Z(\lambda_1-\lambda_0)+\mathbb Z(3\lambda_1)\\
&=\mathbb Z\underbrace{(\lambda_1-\lambda_0)}_{\delta}
 +\mathbb Z\unsure{(1-\rho_M)}\underbrace{(\lambda_1-\lambda_0)}_{\delta}\\
&=\mathbb Z(\delta\unsure{,|)}\oplus\mathbb Z(1-\rho_M)(\delta)\\
&=\mathbb Z\delta+\mathbb Z\rho_M\delta.
\end{aligned}
\]
Note that $M_0$ is stable under $\rho$, \inserted{$\sigma$}, in
general by construction, for $\forall\ x\in M$, we
\[
\begin{aligned}
(109)\qquad x&\equiv\rho_M(x)\quad(M_0)
 &&\text{(i.e. }(1-\rho_M)(x)\in M_0\text{)}\\
x&\equiv\sigma_M(x)\quad(M_0)
 &&(1-\sigma_M)(x)\in M_0.
\end{aligned}
\]
\deleted{and}

It is of course also stable under $\tau_M=-\mathrm{id}_M$.

Thus $M_0$ is \unsure{normalized} by $\mathcal G$, hence
\inserted{as a subgroup of $\mathcal G$)} normal in
$\mathcal G$.
\deleted{It is}\inserted{\uline{It therefore corresponds to}} a normal
subgroup of $\mathcal G_0=\mathrm{Gl}(2,\mathbb Z)$,

\scanpage{64}

containing $[\Pi_0,\Pi_0]$, of index $3$ in $\Pi_0$.
I presume that it \deleted{is} \uline{contains} the
congruence subgroup relative to $12$ --- i.e. that it corresponds to a
\inserted{normal} subgroup of $\mathrm{Sl}(2,\mathbb Z/12\mathbb Z)$.
It is of index $3.12=36$ in $\mathrm{Sl}(2,\mathbb Z)$,
\deleted{$72$ in $\mathrm{Gl}(2,\mathbb Z)$},
\deleted{\unsure{while}} that
\[
\mathrm{Sl}(2,\mathbb Z/12\mathbb Z)\simeq
\underbrace{\mathrm{Sl}(2,\mathbb Z/4\mathbb Z)}_{48}\times
\underbrace{\mathrm{Sl}(2,\mathbb Z/\unsure{3}\mathbb Z)}_{24}
\]
--- it would therefore be a subgroup \uline{\unsure{of order}}
$32$ of $\mathrm{Sl}(2,\mathbb Z/12\mathbb Z)$,
\inserted{it is obviously the normal subgroup}
generated by the elements \deleted{\illegible{1}} the image of
\deleted{$[\sigma_\infty,\lambda_\infty]$},
\[
\left\{
\begin{array}{l}
(\sigma(\rho)\rho^{-1})^3=[\lambda_0,\lambda_1],\ \text{or }[\sigma_\infty,\lambda_\infty]\\
\text{cf. (7)}\\
\delta\overset{\mathrm{dfn}}{=}\lambda_1\lambda_0^{-1}
\ \deleted{\unsure{\in\Pi_0}}\\
\rho(\delta)
\end{array}
\right.
\]
or again the \inserted{normal} subgroup generated by the single element
\deleted{by the images in $\mathrm{Sl}(2,\mathbb Z/12\mathbb Z)$ of)}
\[
\deleted{\delta=\lambda_1\lambda_0^{-1},\quad
\rho(\delta)=\lambda_\infty\lambda_1^{-1},\quad
\rho^2(\delta)=\lambda_0\lambda_\infty^{-1}}
\]
\[
(110)\qquad
\left\{
\begin{aligned}
\delta&=\lambda_1\lambda_0^{-1}=\ell_1\ell_0^{-1}
=\begin{pmatrix}1&0\\2&1\end{pmatrix}
 \begin{pmatrix}1&+2\\0&1\end{pmatrix}\\
&=\begin{pmatrix}1&+2\\2&5\end{pmatrix}
\end{aligned}
\right.
\]
The question is therefore whether this normal subgroup of
$\mathrm{Sl}(2,\mathbb Z/12\mathbb Z)$ is well and truly of order $32$,
hence of index $36$ --- its order could be larger, its index
smaller. To begin with, what is the order of $\delta$?

\scanpage{65}
\mspage{680}

modulo $12$? We have
\[
\begin{array}{ll}
\deleted{\unsure{(111)}}&
\left\{\begin{aligned}
\delta^2&=\begin{pmatrix}1&2\\2&5\end{pmatrix}
\begin{pmatrix}1&2\\2&5\end{pmatrix}
=\begin{pmatrix}5&12\\12&29\end{pmatrix}
\equiv\begin{pmatrix}5&0\\0&5\end{pmatrix}\quad(12)\\
\delta^4&\equiv1\quad(12),
\end{aligned}\right.\\[-1.5em]
\text{whence}&
\end{array}
\]
that is already a good point! We have \deleted{moreover}
\inserted{even}
\[
(112)\qquad
\diagram{$\left\{\begin{array}{lcl}
\delta\equiv1\quad(2)&&\delta\not\equiv1\ (4),\ \delta\not\equiv1\ (3)\\
&\Downarrow&\\
\delta^2\equiv1\quad(4)&&\delta^2\equiv-1\quad(3)
\end{array}\right.$}
\]


Let us consider
\[
(113)\qquad
\diagram{\begin{tikzcd}[column sep=small]
\mathrm{Sl}(2,\mathbb Z/12\mathbb Z)\simeq
\underbrace{\mathrm{Sl}(2,\mathbb Z/4\mathbb Z)}_{48}\times
\underbrace{\mathrm{Sl}(2,\mathbb Z/3\mathbb Z)}_{24}
\arrow[r,"u\times v"]&
\underbrace{\mathcal G_0^+/\Pi_0}_{12}\times
\underbrace{\mathcal G_3^+}_{3}
\end{tikzcd}}
\]
where \unsure{:}
\[
\diagram{\begin{tikzcd}
u:\mathrm{Sl}(2,\mathbb Z/4\mathbb Z)\arrow[r]&\mathcal G_0^+/\Pi_0
\end{tikzcd}}
\]
is the canonical surjective homomorphism, whose kernel is isomorphic to
$\deleted{\unsure{\mathbb Z}}\mu\times\mu$ --- the three nonzero
elements of the kernel are \unsure{just} $\lambda_0,\lambda_1,\lambda_\infty$
permuted among themselves by $\rho$, and
$\dot\lambda_1\dot\lambda_0^{-1}=\dot\lambda_1\dot\lambda_0
=\dot\lambda_\infty$ and $\ldots$ i.e. \deleted{\unsure{the}}
the $\dot\lambda_i$ are none other than $\dot\delta$ and its permuted images
under $\rho$ --- and
\[
\diagram{\begin{tikzcd}
v:\mathrm{Sl}(2,\mathbb Z/3\mathbb Z)\arrow[r]&\mathcal G_3
\end{tikzcd}}
\]
is the composite homomorphism
\[
\diagram{\begin{tikzcd}[row sep=tiny]
\mathrm{Sl}(2,\mathbb Z/3\mathbb Z)\arrow[r,"\text{epimorphism}"]&
\unsure{A_4|G_4}\arrow[r,"\text{epimorphism}"]
\arrow[d,phantom,"\cap" description]&\mathcal G_3^+\\
&\mathcal G_{\mathbb P^1(\mathbb Z/3\mathbb Z)}&
\end{tikzcd}}
\]
whose kernel is of order $8$. The kernel of $u\times v$ is of order
$4.8=32$, and it obviously contains $\dot\delta$. This shows that
it is indeed the \inserted{normal} subgroup generated by $\delta$
(whose order a priori was \deleted{well and} a \uline{multiple} of $32$), and

\scanpage{66}

that \deleted{$\unsure{M_0}$}
\inserted{the subgroup under consideration of $\Pi_0\subset\mathrm{Sl}(2,\mathbb Z)$,
of index $3$ in $\Pi_0$,}
is the inverse image of this preceding subgroup \deleted{in $M$} of
$\mathrm{Sl}(2,\mathbb Z/12\mathbb Z)$.

Let us sum up the results obtained:

\uline{Proposition} Let us consider in $\mathcal G_0^+=\mathrm{Sl}(2,\mathbb Z)$
the tower of invariant subgroups (even in $\mathcal G_0$)
\[
(114)\qquad
\diagram{$\begin{tikzcd}[column sep=small,row sep=small]
\mathcal G_0^+&
\Pi=\Pi[2]\arrow[l,hook',"6"']\arrow[d,equal,"\mathrm{dfn}"]&
\Pi_0\arrow[l,hook',"2"']\arrow[d,equal]&
\Pi_{00}\ (\supset[\Pi_0,\Pi_0])\arrow[l,hook',"3"']\arrow[d,equal]\\
&\begin{array}{c}
\{u\in\mathcal G_0^+\mid u\equiv1\ (2)\}\\
\deleted{\Vert}\\
\deleted{\unsure{\mathfrak{sl}(2,\mathbb F_2)}\simeq
\{u\in\mathrm M_2(\mathbb F_2)\mid\mathrm{Tr}\,u=0\}}\\
\unsure{(i)|(1)}
\end{array}
&\begin{array}{l}
\text{Subgroup}\\
\text{generated}\\
\text{by the }-\ell_i=\lambda_i
\end{array}
&\begin{array}{l}
\text{Subgroup }\uline{\text{invariant}}\\
\text{generated by}\\
\deleted{\text{by the }\lambda_i,}\\
\deleted{[\lambda_0,\lambda_1]\ \text{and}}\\
\lambda_1\lambda_0^{-1}=\ell_1\ell_0^{-1}
\end{array}
\end{tikzcd}$}
\]

Let us consider on the other hand the quotient group
\[
(115)\qquad
\diagram{\begin{tikzcd}[row sep=small,column sep=small]
\mathrm{Sl}(2,\mathbb Z/12\mathbb Z)
\arrow[r,phantom,"\simeq" description]
\arrow[d,phantom,"\simeq" description]&
\mathcal G_0^+/\underbrace{\Pi[12]}_{\simeq\{u\in\mathcal G_0^+\mid u\equiv1\ (12)\}}\\
\mathrm{Sl}(2,\mathbb Z/4\mathbb Z)\times
\mathrm{Sl}(\unsure{2},\mathbb Z/\unsure{3}\mathbb Z),&
\end{tikzcd}}
\]
and the tower of the following subgroups
\[
(116)\qquad
\diagram{$\left\{\begin{array}{rll}
&H=\mathrm{Sl}(2,\mathbb Z/4\mathbb Z)\times\mathrm{Sl}(2,\mathbb Z/3\mathbb Z)
&\text{order }48\times24\\
6\ \cup&\hfill\cup\hfill&\\
&\Pi^\bullet=\mathfrak{sl}(2,\mathbb F_2)\times\mathrm{Sl}(2,\mathbb Z/3\mathbb Z)
&\text{order }8\times24\\
2\ \cup&\hfill\cup\hfill&\\
&\Pi_0^\bullet=\mathfrak{sl}(2,\mathbb F_2)'\times\mathrm{Sl}(2,\mathbb Z/3\mathbb Z)
&\text{order }4\times24\\
3\ \cup&\hfill\cup\hfill&\\
&\Pi_{00}^\bullet=\mathfrak{sl}(2,\mathbb F_2)'\times
\underbrace{\mathrm{Syl}_2}_{\simeq\deleted{\illegible{1}}D_4}
&\text{order }4\times8,
\end{array}\right.$}
\]

\margin{\uline{N.B.} $(\mathcal G_0^+)_{\unsure{\mathrm{ab}}}\simeq
\mathbb Z/12\mathbb Z$ is identified with a
\deleted{$\unsure{(115)|(116)}\ \unsure{\mathrm{Sl}(2)}$}
quotient of $H/\Pi_{00}^\bullet$ by a subgroup (116) of order $3$
--- that is again the subgroup generated by the element
$(-\dot\rho,1)$ of
\[
\underbrace{\mathrm{Sl}(2,\mathbb Z/4\mathbb Z)/\mathfrak{sl}(2,\mathbb F_2)'}_{\mathcal G_0^+/\Pi_0}
\times\mathrm{Sl}(2,\mathbb Z/3\mathbb Z)/\mathrm{Syl}_2.
\]}

Where \deleted{\unsure{we have}}
\[
(117)\qquad
\diagram{$\left\{\begin{array}{ll}
\text{a)}&\mathfrak{sl}(2,\mathbb F_2)\simeq
\mathrm{Ker}\bigl(\mathrm{Sl}(2,\mathbb Z/4\mathbb Z)\longrightarrow
\mathrm{Sl}(2,\mathbb Z/2\mathbb Z)\bigr)\\[.5em]
\text{b)}&\mathfrak{sl}(2,\mathbb F_2)'\unsure{\subset}\mathfrak{sl}(2,\mathbb F_2),
\quad\begin{array}{l}
\text{normal subgroup }\simeq\\
\deleted{\unsure{\mathbb F_3|\mathbb F_2}}\ \mathbb F_2\times\mathbb F_2\ ;\
\text{whose three elements }\ne0\\
\text{are }\dot\lambda_0,\dot\lambda_1,\dot\lambda_\infty\
\text{ (images}\\
\text{in }\mathrm{Sl}(2,\mathbb Z/4\mathbb Z)\text{) --- it is}\\
\text{a complement to its}\\
\text{Lie subalgebra }\mu\simeq\mathbb F_2\ \text{of matrices}\\
\text{\unsure{constant|scalar}}
\end{array}\\[.5em]
\text{c)}&\underset{\substack{\Vert\\
\mathrm{Ker}(\mathrm{Sl}(2,\mathbb Z/3\mathbb Z)\to\mathcal G_3^+)}}{\mathrm{Syl}_2}
\subset\mathrm{Sl}(2,\mathbb Z/3\mathbb Z),\quad
\begin{array}{l}
\text{Sylow subgroup, of order }8\\
\text{(which is normal)}
\end{array}
\end{array}\right.$}
\]

\scanpage{67}
\mspage{681}

This being so, the tower of subgroups (114) is the inverse image of the tower (116).

\deleted{\unsure{Remark|Remark}}
To satisfy my conscience, I am going to take up the preceding \unsure{sketch|sketch}
of a proof again. We note that the homomorphism
$\mathcal G_0^+\longrightarrow H$ maps the groups appearing in
(114) into those of (116) --- for $\Pi$ and $\Pi_0$, we have
in fact known for a while that they are the inverse images of
$\Pi^\bullet$, $\Pi_0^\bullet$! $\unsure{\ast|\leftarrow}$
(For this, it \unsure{was enough|would be enough} to examine
$\mathrm{Sl}(2,\mathbb Z/4\mathbb Z)$, instead of
$\mathrm{Sl}(2,\mathbb Z/12\mathbb Z)$ \ldots).

Thus the inverse image of $\Pi_{00}^\bullet$ in $\mathcal G_0^+$
contains $\Pi_{00}$. As both are of index $36$, these two
groups are equal, q.e.d. (The same
\unsure{argument} could be applied to the case of $\Pi_0$, $\Pi_0^\bullet$ \ldots)

\uline{Corollary} The canonical homomorphism
\[
\diagram{\begin{tikzcd}[column sep=small]
\mathcal G_0^*=\mathrm{Gl}(2,\mathbb Z)\arrow[r]&
\mathrm{Gl}(2,\mathbb Z/12\mathbb Z)/
\bigl(\Pi_{00}^\bullet\deleted{\unsure{\simeq}}=
\mathfrak{sl}(2,\mathbb F_2)'\times\mathrm{Syl}_2\bigr)
\end{tikzcd}}
\]
\unsure{factors} through
\[
\mathcal G=\deleted{\mathcal G_0}\mathrm{Gl}(2,\mathbb Z)/[\Pi_0,\Pi_0]
\]
studied in the present \S\,51, and the inverse image of
$\Pi_{00}^\bullet$ is none other than the subgroup $M_0$ of (99) (108).

\margin{\uline{N.B.}
But $\mathcal G_0^{\unsure{+|*}}\longrightarrow
\unsure{\mathrm{Sl}|\mathrm{Gl}}(2,\mathbb Z/12\mathbb Z)$
does \uline{not} factor through $\mathcal G^{\unsure{+|*}}$,
for if this were so $\Pi_0^{\unsure{\bullet|\prime}}$
would be a \unsure{sub-}quotient of $\Pi_0/[\Pi_0,\Pi_0]$ and
hence would be commutative --- \unsure{now all}
\unsure{we see} \unsure{two|therefore} \unsure{which} \unsure{are}
both.}

\scanpage{68}

The preceding considerations make fairly \unsure{natural|natural} the consideration of the covering of order $3$ of $U_{0,\overline{\mathbb Q}}$ \inserted{\deleted{\unsure{or of} $M_{1,\unsure{1|0}}$}} corresponding to the subgroup $\Pi_{00}$ of $\Pi_0\simeq\Pi_1(U_{0,\overline{\mathbb Q}})$.

\margin{it is even a \unsure{covering} of $M_{1,1}$\dots}

We check at once that it is a curve \uline{of genus $1$} -- if then we take it into our heads to work with $[\Pi_{0,0},\Pi_{0,0}]$ instead of $[\Pi_0,\Pi_0]$, we find a group extension of $\Pi_0/[\Pi_0,\Pi_0]\simeq M$, by a group of rank $2$ which is ``elliptic'', and on which $M$ will operate in a highly nontrivial way -- later on we would have to examine this action, and put it in relation with the action we have \unsure{found}, subordinate to the homomorphism
\diagram{\begin{tikzcd}
\mathrm{Sl}(2,\mathbb Z)^{\wedge}\arrow[r] & \mathrm{Sl}(2,\widehat{\mathbb Z}).
\end{tikzcd}}

\margin{it is an elliptic curve of invariant $0$ (for it is equipped with an \unsure{automorphism} of order $3$)}

\scanpage{69}
\mspage{682}

5) The groups $\underline M_{\rho,\sigma}$, $\underline M_{\rho,\sigma}^{\sharp}$ (via $\underline N_{\rho,\sigma}$).

To every $\ell\in\Gamma\unsure{\underline M_0}$ there corresponds an inner automorphism $\operatorname{int}(\ell)$, which is the identity on $\underline M_0$. We will have, tautologically
\[
(118)\qquad
\left\{\begin{aligned}
\operatorname{int}(\ell)(\rho)&=\alpha_{\unsure{0|\cdot}}(\rho)\\
\operatorname{int}(\ell)(\sigma)&=\beta(\sigma)
\end{aligned}\right.
\qquad\alpha=\beta=\ell
\]

But \deleted{we will have} will we for all that have the loop relation (relative to $\underline\mu=1$ hence $\nu''=0$)
\[
(1-\rho_M)(\underbrace{\alpha}_{\ell})=(1-\sigma_M)(\underbrace{\beta}_{\ell}).
\]
\[
\underset{\unsure{(*)|\alpha_1}}{\text{i.e.}}\qquad
\rho_M(\ell)=\sigma_M(\ell)\qquad ?
\]

Setting again
\[
(119)\qquad\ell=c_0\lambda_0+c_1\lambda_1
\]
\deleted{we will have}
the relation $(*)$ is written
\[
c_0\lambda_1+c_1\underbrace{\lambda_\infty}_{-\lambda_0-\lambda_1}
=c_0\lambda_1+c_1\lambda_0
\]
\[
\text{i.e.}\qquad-\unsure{c_1}\lambda_0+(c_0-c_1)\lambda_1
=c_1\lambda_0+c_0\lambda_1
\]
i.e. $c_1=0$ i.e. $\ell=c_0\lambda_0$. Thus it is subject to the \unsure{condition|conditions} that we restrict ourselves to considering $\ell\in\Gamma\underline L_0^{\unsure{\underline M_0|\underline M}}\unsure{\Rightarrow}$; \unsure{corresponding|corresponding}

\diagram{\begin{tikzcd}[column sep=small]
{(120)\quad\deleted{\unsure{\underline L_0}}\quad\ell=c_0\lambda_0\in
\underline L_{0\deleted{\unsure{,\rho,\sigma}}}^{\underline M}
\overset{\mathrm{dfn}}{=}\operatorname{Im}(}
& \underline{\mathbb Z}\arrow[r]\arrow[d,phantom,"\simeq" description]
& \underline M)\\
& \mathbb G_a\arrow[d,phantom,"\unsure{\ni}" description] &\\
& c_0\arrow[r,mapsto] & c_0\lambda_0
\end{tikzcd}}

We have a \deleted{\unsure{dia}} diagram
\diagram{\begin{tikzcd}[column sep=large]
{(121)\quad\mathbb G_a\simeq\underline L_0^{\unsure{\underline M_0|\underline M}}}
\arrow[r]\arrow[d]
& \underline N_{\rho,\sigma} & \deleted{\unsure{\underline L_0^{\underline M}}}\\
\underline{\mathcal G}^{+} & &
\end{tikzcd}}

\scanpage{70}

the image of
\diagram{\begin{tikzcd}
{(\unsure{127|122})\quad\underline L_0^{\unsure{\underline M_0|\underline M}}}\arrow[r]
& \underline N_{\rho,\sigma}
\end{tikzcd}}
is invariant -- the homomorphism $(121)$ being in fact compatible with the action of $N_{\rho,\sigma}$ -- via $\underline\mu$ on $\underline L_0^{M_0}\simeq\mathbb G_a$ (i.e. induced by the operation of $\underline N_{\rho,\sigma}$ \unsure{of} on\unsure{!| the} $\underline{\mathcal G}^{+}\supset\underline M\supset\underline L_0^{\underline M}$), and \deleted{the}\inserted{the two} operations of \deleted{$\unsure{N_\rho|A_\rho}$} $\underline L_0^{\underline M}$ on $\underline{\mathcal G}^{+}$ deduced from $(121)$ being \unsure{identical|identical}. Thus according to the general \unsure{theory|theorem} we \unsure{find|will find} a group scheme $\underline M_{\rho,\sigma}$ and a diagram of exact sequences

\diagram{\begin{tikzcd}[column sep=large,row sep=large]
{(123)\quad1}\arrow[r]
&\underline L_0^{\underline M}\arrow[r]\arrow[d]
&\underline N_{\rho,\sigma}\arrow[r,Rightarrow]\arrow[d]
&\overset{\inserted{\overset{\mathrm{dfn}}{=}\underline N_{\rho,\sigma}/\underline L_0^{\underline M}}}{\underline\Gamma_{\rho,\sigma}}\arrow[r]\arrow[d,Rightarrow]
&1\\
1\arrow[r]
&\underline{\mathcal G}^{+\unsure{a|\alpha}}\arrow[r]
&\underline M_{\rho,\sigma}\arrow[r,Rightarrow]
&\underline\Gamma_{\rho,\sigma}\arrow[r]
&1
\end{tikzcd}}

\margin{with the notations of \S\,49 $\overline{\underline{\mathrm{VI}}}$, the group denoted here $\underline M_{\rho,\sigma}$ would probably be denoted $\underline M_{\rho,\sigma}/\unsure{H_\rho|H_{\underline\Pi}}$, and $\underline M_{\rho,\sigma}^{\sharp}$ would be $\underline M_{\rho,\sigma}^{\sharp}/\unsure{H_\rho}$ (loc. cit. $8^\circ$)}

Replacing $\underline{\mathcal G}^{+}$ by $\underline M$, we likewise find $\underline M_{\rho,\sigma}^{\sharp}$, and

\diagram{\begin{tikzcd}[column sep=large,row sep=large]
{(124)}\arrow[d,phantom,"\Bigg\{" description]
&1\arrow[r,dash]
&\underline L_0^{\underline M}\arrow[r,dash]\arrow[d]
&\underline N_{\rho,\sigma}\arrow[r,dash]\arrow[d]
&\underline\Gamma_{\rho,\sigma}\arrow[r,dash]\arrow[d,Rightarrow]
&1\\
{} &1\arrow[r]
&\deleted{\unsure{\underline{\mathcal G}^{+}}}\underline M\arrow[r]
&\underline M_{\rho,\sigma}^{\sharp}\arrow[r]
&\underline\Gamma_{\rho,\sigma}\arrow[r]
&1
\end{tikzcd}}
,

moreover we have a homomorphism from $(124)$ into $(123)$ (compare \S\,\uline{49} $\overline{\underline{\mathrm{VI}}}$ $5^\circ$\dots), and
\[
(125)\qquad M_{\rho,\sigma}^{\sharp}\subset M_{\rho,\sigma}\qquad\dots
\]

\scanpage{71}
\mspage{683}

We must specify the structure of $\Gamma_{\rho,\sigma}$, noting that $L_0^{\underline M}\subset N_{\rho,\sigma}^{0+}$, hence
\diagram{\begin{tikzcd}
{(\unsure{126|125})\quad1}\arrow[r]
&\deleted{\unsure{N_{\rho,\sigma}}}\underline\Gamma_{\rho,\sigma}^{0+}\arrow[r]
&\underline\Gamma_{\rho,\sigma}\arrow[r]
&\underline{\mathbb O}_{12}^{*}\arrow[r]
&1
\end{tikzcd}}
with
\diagram{\begin{tikzcd}
{(127)\quad1}\arrow[r]
&\underline L_0^{M}\arrow[r,thick]
&\underline N_{\rho,\sigma}^{0+}\arrow[r]
&\underline\Gamma_{\rho,\sigma}^{0+}\arrow[r]
&1\\
&&\deleted{\unsure{\text{it is }C_f|\text{it is }G_f}}&&
\end{tikzcd}}

Let us consider \deleted{the}\inserted{the}\deleted{exact sequence}\inserted{surjective homomorphism} of $\mathbb Z$-modules

\diagram{\begin{tikzcd}[column sep=small,row sep=small]
{(128)\quad\unsure{N_{0,\sigma}^{0+}|N_{\rho,\sigma}^{0+}}=
\deleted{\unsure{M\times M}}\operatorname{Ker}(}
&M\times M\arrow[r]
&M_0)\deleted{\unsure{\longrightarrow0}}\arrow[d,phantom,"\cap" description]\\
&&M
\end{tikzcd}}
\diagram{\begin{tikzcd}[column sep=small]
\deleted{\unsure{N}}\arrow[r,mapsto]
&\deleted{\unsure{\alpha_0,\beta_0}}\;\deleted{\unsure{(\alpha_0,\beta_0)}}
\end{tikzcd}}
\diagram{\begin{tikzcd}[column sep=large]
{(\underset{\substack{\Vert\\\inserted{a_0\lambda_0+a_1\lambda_1}}}{\alpha_0},
\underset{\substack{\Vert\\\inserted{b_0\lambda_0+b_1\lambda_1}}}{\beta_0})}\arrow[r]
& (1-\rho_M)(\alpha_0)-(1-\sigma_M)(\beta_0)
\end{tikzcd}}
\[
\underbrace{\hphantom{(1-\rho_M)(\alpha_0)-(1-\sigma_M)(\beta_0)}}
\]
\inserted{kernel consisting of the $(\alpha_0,\beta_0)$ such that
\[
\begin{aligned}
c_0&\overset{\mathrm{dfn}}{=}a_0+a_1+(b_1-b_0)=0\\
c_1&\overset{\mathrm{dfn}}{=}2a_1-a_0-(b_1-b_0)=0
\end{aligned}
\qquad\bigg|\ \text{cf. }(61)
\]}

Let $\underline N^{0+}$ be the kernel, so that we have
\[
(129)\qquad\underline N_{\rho,\sigma}^{0+}\simeq\underline N^{0+}.
\]

\uline{Let us consider} the homomorphism
\diagram{\begin{tikzcd}
{(130)}\arrow[d,phantom,"\Bigg\{" description]
&\mathbb Z\arrow[r]&N^{0+}\\
{} &n\arrow[r,mapsto]&(n\lambda_0,n\lambda_0)
\end{tikzcd}}

the image \deleted{is} is interpreted as \unsure{made up|made up} of the systems $(a_0,a_1,b_0,b_1)$ such that \deleted{$\unsure{a_0=b_0}$}
\[
a_1=b_1=0,\qquad a_0=b_0,\qquad\unsure{''}
\]
\deleted{whence therefore \unsure{a} \illegible{2}}
\inserted{i.e. we have an exact sequence}
\diagram{\begin{tikzcd}
\deleted{\unsure{(131)}}&
0\arrow[r]&\mathbb Z\arrow[r]&N^{0+}\arrow[r,thick]&\mathbb Z^3\\
&&&(a_0,a_1,b_0,b_1)\arrow[r,mapsto]&(a_1,b_1,a_0-b_0)
\end{tikzcd}}
hence \deleted{\unsure{this}} $N^{0+}/\mathbb Z$ is \deleted{\unsure{flat|free}} a free $\mathbb Z$-module \inserted{of rank $1$},

\scanpage{72}

thus \unsure{an} exact sequence of $\mathbb Z$-modules
\diagram{\begin{tikzcd}[column sep=small,row sep=small]
{(131)\quad1}\arrow[r,dash]
&\mathbb Z\arrow[r]
&N^{0+}\arrow[r]
&\Gamma^{0+}\arrow[r]\arrow[d,phantom,"{\rotatebox{90}{$\simeq$}\ \mathrm{dfn}}" description]
&0\\
&n\arrow[r,mapsto]&(n\lambda_0,n\lambda_0)
&N^{0+}/\operatorname{Im}\mathbb Z&
\end{tikzcd}}
with $\Gamma^{0+}$ a free $\mathbb Z$-module of rank $1$.

We conclude from this
\[
(132)\qquad\Gamma_{\rho,\sigma}^{0+}\simeq\underline\Gamma^{0+}
\qquad(\simeq\mathbb G_a)
\]

\deleted{\unsure{Hence}} Thus, all the group schemes we have just constructed are smooth over $\operatorname{Spec}\mathbb Z$.

Can we define a more-or-less natural section of the extension
(\inserted{$N_{\rho,\sigma}$ of $\Gamma_{\rho,\sigma}$ by $L_0^{\underline M}$, hence of}
$M_{\rho,\sigma}^{\sharp}$ of
$\underbrace{\Gamma_{\rho,\sigma}}_{\inserted{\substack{\text{smooth of relative}\\\text{dimension }2}}}$
by $\underline M$?

Taking into account the loop equation which describes $N_{\rho,\sigma}$, it seems that we \unsure{find} such a section \unsure{as soon|some} as we choose a splitting of
\diagram{\begin{tikzcd}[column sep=small,row sep=large]
{} &0\arrow[r]
&\overset{\inserted{\mathrm{rg}\,2}}{N^{0+}}\arrow[r]
&M\times M\arrow[r]
&\overset{\inserted{\mathrm{rg}\,2}}{M_0}\arrow[r]
&0&\text{and of}\\
{(133)}\arrow[u,phantom,"\Bigg\{" description]
&0\arrow[r]
&\underset{\inserted{\mathrm{rg}\,1}}{\mathbb Z}\arrow[r]
&N^{0+}\arrow[r]
&\underset{\inserted{\mathrm{rg}\,1}}{\Gamma^{0+}}\arrow[r]
&0&
\end{tikzcd}}
(category of $\mathbb Z$-modules). But we must check that the \textquotedblleft set-theoretic\textquotedblright{} sections of schemes that \unsure{this} allows us to define are multiplicative\dots

\scanpage{73}
\mspage{684}

\(6^\circ)\) \(\underline Z_{\rho,\sigma}\) and
\(\underline G_{\rho,\sigma}=\mathbb O_{12}^{*}\mathbin{.}\underline{\mathcal G}^{+},\quad
\underline G_{\rho,\sigma}^{\sharp}\simeq\mathbb O_{12}^{*}\mathbin{.}\underline{\mathcal G}^{\sharp}
=\mathbb O_{12}^{*}\mathbin{.}\underline M\).

Once \(\mathbb O_{12}^{*}\) and \(M_0\subset M\) have been introduced, it
seems to me that the group \(\underline Z_{\rho,\sigma}\) (in place
of \(\underline{\widetilde Z}_{\rho,\sigma}\), in which we took
account of the \(\alpha,\beta\) and not only
of \(\xi,\eta\)) \unsure{again becomes} nice:
\[
(134)\quad
\begin{gathered}
\underline Z_{\rho,\sigma}\subset
\underline{\operatorname{Aut}}(\underline{\mathcal G}^{+})\times\mathbb O_{12}^{*}
\\[-2pt]
\text{\(\prime\prime\)}\\[-2pt]
\left\{(u,\underline\mu)\ \left|\
\begin{array}{ll}
\text{a)}&u(\rho)=\xi\rho\quad\text{(locally) }\underline{\mathcal G}_{\tau}^{\sharp}\text{-conjugate to }\rho,\text{ by an}\\
&\alpha\in\underline{\mathcal G}_{\tau}^{\sharp}\text{ such that }\unsure{\chi}(\alpha)=\varepsilon_3(\underline\mu)\\
\text{b)}&u(\sigma)=\eta\sigma\quad\text{(locally) }\underline{\mathcal G}_{\tau}^{\sharp}\unsure{+|-}\text{conjugate to }\sigma,\text{ by a}\\
&\beta\in\underline{\mathcal G}_{\tau}^{\sharp}\text{ such that }\chi(\beta)=\varepsilon_2(\underline\mu)\\
\text{c)}&\text{loop equation}\\
&\xi_0=\eta_0(\lambda_{\unsure{1}}^3)^{\nu''}\quad\text{i.e.}\quad\xi_0=\eta_0+3\nu''\lambda_1\\
&\hfill\text{in }\underline M_0
\end{array}\right.\right.
\end{gathered}
\]
where \(\nu''=\nu''(\underline\mu)\in\Gamma\mathbb G_a\)
(cf. \unsure{page\ verso} page~686), and \unsure{above\(^{t}\)})
\[
(135)\qquad
\left\{\begin{aligned}
\xi_0&=\xi-\nu_3\lambda_1\in\Gamma\underline M_0
&\nu_3&=\nu_3(\underline\mu)\\
\eta_0&=\eta\omega^{\nu_2}\in\Gamma\underline M_0
&\nu_2&=\nu_2(\underline\mu)
\end{aligned}\right.
\]
\deleted{We have that} \inserted{We shall have} \deleted{must suppose} moreover that
\[
(136)\qquad
\left\{\begin{aligned}
&\xi_0\quad\text{\unsure{locally|locally\(^{mt}\)} in the image of }(1-\rho_{\underline M})(\underline M)\\
&\eta_0\quad\text{\unsure{locally|locally\(^{mt}\)} in the image }(1-\sigma_{\underline M})(\underline M)
\end{aligned}\right.
\]
and we shall have
\[
(137)\quad
\begin{aligned}
\underline Z_{\rho,\sigma}\simeq\underline\Sigma_{\rho,\sigma}^{\mathrm{eff}}
\overset{\mathrm{dfn}}{=}
\bigl\{(\underline\mu,\xi_0,\eta_0)&\in\Gamma\mathbb O_{12}^{*}\times\underline M_0\times M_0\ \big|\\
1^\circ)\quad\xi_0&=\eta_0+\deleted{\unsure{\xi}}\nu''(3\lambda_1)\quad\text{in }\underline M_0\\
&\text{where }\nu''=\nu''(\underline\mu)\\
2^\circ)\quad&\text{we have }(136)\quad\deleted{\unsure{\bigr\}}}\ \bigr\}.
\end{aligned}
\]
Now

\((1-\rho_{\underline M}):\)
\diagram{\begin{tikzcd}[column sep=large]
\underline M\arrow[r]&\underline M_0
\end{tikzcd}}

\((1-\sigma_{\underline M}):\)
\diagram{\begin{tikzcd}[column sep=large]
\underline M\arrow[r]&\underline{\mathbb Z}=\mathbb G_a\arrow[r]&\underline M_0\\
(\unsure{b_0}\lambda_0+b_1\lambda_1)\arrow[r]&(b_1-\underset{\deleted{\unsure{4|\lambda}}}{b_0})\,\deleted{\unsure{\lambda}}\arrow[r,mapsto]&\lambda\underbrace{(\lambda_1-\lambda_0)}_{\delta}=\lambda\delta
\end{tikzcd}}

\scanpage{74}

thus (136) \deleted{\unsure{is written}} is a condition on
\(\eta_0\) alone, namely

\((136\mathrm{bis})\quad\eta_0\in\Gamma\operatorname{Im}\bigl(\)
\diagram{\begin{tikzcd}[column sep=large]
\mathbb G_a\arrow[r]&\underline M_0\\
\lambda\arrow[r,mapsto]&\lambda\delta
\end{tikzcd}}
\(\bigr)=\unsure{\Gamma}\operatorname{Ker}\bigl(\)
\diagram{\begin{tikzcd}[column sep=large]
M_0\arrow[r]&\mathbb G_a\\
c_0\lambda_0+c_1\lambda_1\arrow[r]&c_0+c_1
\end{tikzcd}}
\(\bigr)\)

\deleted{That said, we thus find \unsure{as}}
\inserted{Let us note that we have}
\[
\deleted{(138)\qquad\underline\Sigma_{\rho,\sigma}^{\mathrm{eff}}\simeq}
\]
\[
(138)\qquad
\left\{\begin{aligned}
M_0&=\mathbb Z\underbrace{\delta}_{\lambda_1-\lambda_0}
\oplus\mathbb Z\mathbin{.}\underbrace{(3\lambda_1)}_{(1-\rho)\delta}
\qquad\text{whence}\\
\underline M_0&=\deleted{\unsure{\mathbb G_a}}\ \underline{\mathbb Z\delta}
\times\underline{\mathbb Z\mathbin{.}(3\lambda_1)}\ .
\end{aligned}\right.
\]
and \deleted{consequently} we find

\((139)\)
\diagram{\begin{tikzcd}[column sep=large]
\underline\Sigma_{\rho,\sigma}^{\mathrm{eff}}\arrow[r,"\simeq"]&
\mathbb O_{12}^{*}\times\deleted{\unsure{\underline{\mathbb Z}\delta}}\ \deleted{\unsure{\mathbb G_a}}\ \underset{\unsure{\in|\simeq}\mathbb G_a}{\underline{\mathbb Z\delta}}\\
(\underline\mu,\xi_0,\unsure{\eta_0|\beta_0})\arrow[r,mapsto]&(\underline\mu,\unsure{\eta_0|\beta_0})
\end{tikzcd}}
.

For \(\unsure{\xi_0|\zeta_0}\) is recovered with \(\underline\mu,\unsure{\eta_0|\beta_0}\)
\deleted{\unsure{from}} the loop equation
\[
\deleted{\boxed{\begin{aligned}
\unsure{\alpha_0|\xi_0}&=(1-\rho_{\underline M})^{-1}
\bigl(\unsure{\eta_0|\beta_0}+\nu''\underbrace{(3\lambda_1)}_{(1-\rho_{\underline M})(\delta)}\bigr)\\
&\deleted{\unsure{=\nu''\delta+(1-\rho_{\underline M})^{-1}\eta_0}}\\
&=\nu''\delta+\unsure{\xi_0|\eta_0}
\end{aligned}}}
\]

\uline{N.B.} If on the other hand we give ourselves \(\xi_0\), we obtain from it
by the loop equation and (138) the value of \(\unsure{\eta_0|\beta_0}\),
and that of \(\nu''\) -- thus if we know
in addition \(\nu_2\) and \(\nu_3\), we deduce
the value of \(2\nu_0\) hence (over a
base where 2 is a regular section of \(\mathscr O_S\)) \(\nu_0\),
i.e. finally \(\underline\mu\). Thus over a

\scanpage{75}
\mspage{685}

base \(S\) where 2 is regular, we find

\((140)\)
\diagram{\begin{tikzcd}[column sep=large]
\underline\Sigma_{\rho,\sigma}^{\mathrm{eff}}(S)\arrow[r,hook]&
\bigl(\underline M_0\deleted{\unsure{(S)}}\times\{0,1\}\deleted{\unsure{(S)}}\times\{0,1\}\bigr)(S)\\
(\underline\mu,\xi_0,\eta_0)\arrow[r,mapsto]&(\xi_0,\varepsilon_2(\underline\mu),\varepsilon_3(\underline\mu))
\end{tikzcd}}

the image being formed of the systems
\((\xi_0,\unsure{\varepsilon_2|\xi_2},\unsure{\varepsilon_3|\xi_3})\) such that, setting
\[
\deleted{(\unsure{141|140})}\qquad
\xi_0=\unsure{\eta_0|\beta_0}\delta+\nu''(3\lambda_1)
\qquad\eta_0,\nu''\in\deleted{\unsure{\Gamma}}\mathbb G_a(S)
\]
thanks to (138), we have
\[
(141)\qquad
\left\{\begin{aligned}
&\nu''+\nu_2-\nu_3[+2\nu_2\nu_3]\quad
\begin{gathered}
\inserted{\text{divisible by 2}}\\[-2pt]
\deleted{\text{of the form }2\nu_0}\\[-2pt]
\deleted{\text{divisible by}}
\end{gathered}\\
&\underbrace{(-5)^{\nu_2}(+5)^{\nu_3}+6(\nu''+\nu_2-\nu_3+2\nu_2\nu_3)}_{1-2\nu_3-2\nu_2\nu_3+6\nu''}\quad\text{inverse}
\end{aligned}\right.
\]

In any case we see that
\(\underline\Sigma_{\rho,\sigma}^{\mathrm{eff}}\) hence also \(\underline Z_{\rho,\sigma}\) is smooth of
relative dimension 2

\((142)\)
\diagram{\begin{tikzcd}[column sep=large]
1\arrow[r,dash]&\unsure{\underline Z}_{\rho,\sigma}^{+0}\arrow[r]\arrow[d,phantom,"\rotatebox{90}{$\simeq$}"']&
\underline Z_{\rho,\sigma}\arrow[r]&\mathbb O_{12}^{\unsure{*|\circ}}\arrow[r,dash]&1\\
&\mathbb G_a\simeq\underline{\mathbb Z\delta}
\end{tikzcd}}

(the isomorphism
\diagram{\begin{tikzcd}[column sep=large,row sep=small]
&&\inserted{\underline\Sigma_{\rho,\sigma}^{\mathrm{eff}+0}}\\
\mathbb G_a\arrow[r,phantom,"\simeq"]&\underline Z_{\rho,\sigma}^{+0}\arrow[ur,"\inserted{\simeq}"]&
\deleted{\unsure{\mathbb G_a}}\arrow[l,phantom,"\deleted{\unsure{\xleftarrow{\simeq}}}"']
\end{tikzcd}}
being given by
\diagram{\begin{tikzcd}[column sep=large]
\lambda\arrow[r,mapsto]&(\underline\mu=1,\xi_0\deleted{\unsure{=\eta_0}}=\lambda\delta,\eta_0=\xi_0=\lambda\delta)
\end{tikzcd}}
)

We \deleted{shall have} shall set
\[
(143)\qquad
\unsure{\underline L_0^{\sharp}|L_0^{\sharp}}\overset{\mathrm{dfn}}{=}
\{\lambda\mathbin{.}\lambda_0\mid\lambda\in\Gamma\mathbb G_a\}\subset\underline M
\]
and we define

\((144)\)
\diagram{\begin{tikzcd}[column sep=large]
\underline L_0^{\sharp}\arrow[r]&\underline Z_{\rho,\sigma}^{\unsure{+0}}\arrow[r,"\simeq","\bigl("{description,pos=0,anchor=east}]&
\underline\Sigma_{\rho,\sigma}^{\mathrm{eff}\unsure{+0}}\bigr)\\
\underset{\prime\prime\lambda\mathbin{.}\lambda_0}{\ell}\arrow[r,mapsto]&
(\operatorname{int}(\ell),1)\arrow[r,mapsto]&
\begin{gathered}
(\deleted{\unsure{\mu}}1,(1-\rho_{\underline M})(\ell),(1-\sigma_{\underline M})(\ell))\\
=(1,-\lambda\delta,-\lambda\delta)
\end{gathered}
\end{tikzcd}}

\scanpage{76}

thus we have
\[
(145)\qquad\underline L_0^{\sharp}\simeq\underline Z_{\rho,\sigma}^{+0}
\]
we shall have
\[
(146)\qquad\underline\gamma_{\rho,\sigma}\overset{\mathrm{dfn}}{=}
\underline Z_{\rho,\sigma}/\underline L_0^{\sharp}\simeq\underline{\mathbb O}_{12}^{*}\ .
\]
If we now construct \(G_{\rho,\sigma}\)

\deleted{\((\unsure{147})\)}

as
\[
(147)\qquad\underline G_{\rho,\sigma}
=(\underline Z_{\rho,\sigma}\mathbin{.}\underline{\mathcal G}^{+})/\underline L_0^{\sharp}
\]
we find

\((148)\)
\diagram{\begin{tikzcd}[column sep=large,row sep=large]
1\arrow[r]&\underline L_0^{\sharp}\arrow[r]\arrow[d]&
\underline Z_{\rho,\sigma}\arrow[r]\arrow[d]&
\overbrace{\gamma_{\rho,\sigma}}^{\mathbb O_{12}^{*}}\arrow[r]\arrow[d,phantom,"\rotatebox{90}{$\simeq$}" description]&1\\
1\arrow[r,dash]&\underline{\mathcal G}^{+}\arrow[r]&G_{\rho,\sigma}\arrow[r]&\gamma_{\rho,\sigma}\arrow[r]&1
\end{tikzcd}}
.

Moreover we have a canonical splitting
of \(\underline Z_{\rho,\sigma}\) over \(\gamma_{\rho,\sigma}=\mathbb O_{12}^{*}\)

\((149)\)
\diagram{\begin{tikzcd}[column sep=large]
\mathbb O_{12}^{\unsure{*|\wedge}}\arrow[r]&\underline Z_{\rho,\sigma}\arrow[r,"\simeq"]&\underline\Sigma_{\rho,\sigma}^{\mathrm{eff}}\\
\underline\mu\arrow[r,mapsto]&\unsure{u_{\underline\mu}|\underline\mu}\arrow[r]&
(\underline\mu,\xi_0=\nu''(3\lambda_1),\eta_0=0)
\end{tikzcd}}

\margin{\uline{N.B.} We set
\(G_{\rho,\sigma}^{\sharp}=\underbrace{\mathbb O_{12}^{*}}_{\gamma_{\rho,\sigma}}
\mathbin{.}\underset{\underline M}{\underline{\mathcal G}^{\sharp}}\subset G_{\rho,\sigma}\).}

whence
\[
(150)\qquad
\left\{\begin{aligned}
\underline Z_{\rho,\sigma}&\simeq\mathbb O_{12}^{*}\mathbin{.}\underline L_0^{\sharp}\\
G_{\rho,\sigma}&\simeq\mathbb O_{12}^{*}\mathbin{.}\underline{\mathcal G}^{+}
\end{aligned}\right.
\]
The operation of \(\mathbb O_{12}^{*}\) on \(\underline L_0^{\sharp}\) is
via
\diagram{\begin{tikzcd}[column sep=large]
\mathbb O_{12}^{*}\arrow[r,dash]&\mathbb G_m
\end{tikzcd}}
and homotheties, that
of \(\mathbb O_{12}^{*}\) on \(\underline{\mathcal G}^{+}\) is given by
\[
(151)\quad
\left\{\begin{aligned}
u_{\underline\mu}(\rho)&=\bigl[\deleted{\unsure{3\nu''|\nu_3}}\unsure{6\nu''|6\nu}\lambda_1\bigr]\rho
\quad\text{(}=\lambda_1^{\unsure{6\nu''|6\nu}}\rho\quad\text{in \deleted{more} exponential notation}\text{)}\\
u_{\underline\mu}(\sigma)&=\unsure{\omega^{\nu_2}}\sigma\\
\deleted{\unsure{u_{\underline\mu}}}\qquad u_{\underline\mu}(x)&=\mu x\qquad\text{if }x\in\Gamma\underline M\ .
\end{aligned}\right.
\]
